% Stronger unconditional version retaining the flight-graph mixed volume.
% Intended to replace the all-loose bound after referee-style audit.

\subsection{Relative regularization and an unconditional graph-structured bound}
\label{subsec:relative-graph-bound}

Fix a cyclic $k$-flight itinerary.  Write $q_j=e^{it_j}$ and let
$\tau_j\in(0,2\pi)$ be the positively oriented gap from $q_j$ to $q_{j+1}$,
with $q_{k+1}=q_1$ and $\sum_j\tau_j=2\pi$.  Introduce one positive private
variable per flight,
\[
 y_j=\begin{cases}
 R_j=e^{a_j\tau_j},&a_j\ne0,\\
 T_j=\tau_j,&a_j=0.
 \end{cases}
\]
Let $\mathcal C_k$ be the corresponding connected cyclic configuration chamber
and set
\[
 \mathcal M_k=\mathcal C_k\times(0,\infty)^k.
\]
For the matching equations use the local boxes
\begin{equation}\label{eq:rg-local-box}
 L_j=[-M,M]_{q_j}+[-M,M]_{q_{j+1}}+[0,1]_{y_j}.
\end{equation}
These contain the Newton supports of the nonzero-slope Laurent flight relation
and of the zero-slope polynomial--Laurent relation described above.

Let $\phi_j=\log y_j-a_j\tau_j$ on a nonzero-slope flight and
$\phi_j=y_j-\tau_j$ on a zero-slope flight, and define the partial physical
leaves
\begin{equation}\label{eq:rg-Gamma-r}
 \Gamma_r=\{\phi_1=\cdots=\phi_r=0\}\subset\mathcal M_k,
 \qquad r=0,\ldots,k,
\end{equation}
with $\Gamma_0=\mathcal M_k$.  Thus $\Gamma_k$ is the physical flight leaf.
Let
\[
 Q=(Q_1,\ldots,Q_k):\mathcal M_k\longrightarrow\mathbb R^k
\]
be the Laurent matching map.

\begin{lemma}[Separating flight chain and uniform covering]
\label{lem:rg-covering}
The partial leaves \(\Gamma_r\) form a separating chain on \(\mathcal M_k\).
Moreover
\begin{equation}\label{eq:rg-rho}
 \rho_k(q,y)=\prod_{j=1}^k
 \left(2-\frac{q_{j+1}}{q_j}-\frac{q_j}{q_{j+1}}\right)
 \frac{y_j}{(1+y_j)^2}
\end{equation}
is a covering function.  For every sufficiently small regular
\(\varepsilon>0\), the same real level is represented after clearing positive
denominators by
\begin{equation}\label{eq:rg-Peps}
 P_\varepsilon=
 \prod_{j=1}^k
 \left(2-\frac{q_{j+1}}{q_j}-\frac{q_j}{q_{j+1}}\right)y_j
 -\varepsilon\prod_{j=1}^k(1+y_j)^2=0.
\end{equation}
\end{lemma}

\begin{proof}
Each \(d\phi_j\) has a nonzero component in the private \(y_j\)-direction,
so \(\Gamma_r\) is a smooth cooriented hypersurface in \(\Gamma_{r-1}\) and
is the boundary there of \(\{\phi_j<0\}\cap\Gamma_{r-1}\).
The phase factor in \eqref{eq:rg-rho} equals
\[
 4\sin^2(\tau_j/2)>0.
\]
The closure of the cyclic configuration chamber in \((S^1)^k\) is compact.
Thus a sequence leaving every compact subset of \(\mathcal M_k\) either has
some collision \(\tau_j\to0\), or some \(y_j\to0\) or \(y_j\to\infty\).
In all cases \(\rho_k\to0\), which is exactly the covering property.  Sard's
theorem supplies arbitrarily small regular values.  Multiplication by
\(\prod_j(1+y_j)^2>0\) gives \eqref{eq:rg-Peps} without changing the real
level or its coorientation.
\end{proof}

\begin{lemma}[Simultaneous rigid-level regularization]
\label{lem:rg-regularization}
Let $Z\subset\Gamma_k\cap Q^{-1}(0)$ be a finite collection of points at which
$d(Q|_{\Gamma_k})$ has full rank $k$.  Then there are arbitrarily small
$b\in\mathbb R^k$ such that:
\begin{enumerate}
\item $b$ is a regular value of $Q|_{\Gamma_r}$ for every $r=0,\ldots,k$;
\item every point of $Z$ has a unique continuation to a point of
$\Gamma_k\cap Q^{-1}(b)$ in a prescribed sufficiently small neighborhood;
\item on the regular level $V_b=Q^{-1}(b)$ the chain
\[
 V_b\cap\Gamma_k\subset\cdots\subset V_b\cap\Gamma_1\subset V_b
\]
is a separating chain for the restricted forms $d\phi_1,\ldots,d\phi_k$.
\end{enumerate}
The shifts $Q_j\mapsto Q_j-b_j$ do not enlarge the local Newton boxes $L_j$.
\end{lemma}

\begin{proof}
At every point of $Z$, the map $Q|_{\Gamma_k}$ is a local diffeomorphism.
Choose pairwise disjoint neighborhoods of the points of $Z$ and a common ball
$B_\delta(0)\subset\mathbb R^k$ such that every $b\in B_\delta(0)$ has one
continued preimage in each neighborhood.

For each $r=0,\ldots,k$, Sard's theorem implies that the critical values of
$Q|_{\Gamma_r}$ have measure zero.  Because there are only finitely many
$r$, choose $b\in B_\delta(0)$ outside the union of these critical-value sets.
This proves the first two statements.

It remains to verify the separating chain on $V_b$.  Regularity of
$Q|_{\Gamma_{r-1}}$ makes $V_b\cap\Gamma_{r-1}$ smooth.  If
$d\phi_r$ vanished on the tangent bundle of this level at a point of
$V_b\cap\Gamma_r$, then $d\phi_r$ would lie in the span of the components of
$dQ$ on $T\Gamma_{r-1}$.  Restricting $dQ$ to
$T\Gamma_r=\ker(d\phi_r|_{T\Gamma_{r-1}})$ would then lower its rank below
$k$, contradicting that $b$ is a regular value of $Q|_{\Gamma_r}$.  Hence
$\phi_r=0$ is a regular hypersurface in $V_b\cap\Gamma_{r-1}$.  Since
$V_b\cap\Gamma_r$ is the boundary there of
$\{\phi_r<0\}\cap V_b\cap\Gamma_{r-1}$ with the coorientation induced by
$d\phi_r$, it is separating.

Finally, the origin belongs to every box $L_j$, so subtraction of the constant
$b_j$ creates no exponent outside $L_j$.
\end{proof}

We next isolate the relative form of the Khovanskii--Rolle argument that is
needed here.  It is the same one-dimensional Rolle step used in the standard
proof of the Khovanskii reduction, but applied on the regular algebraic level
$V_b$.

\begin{proposition}[Relative Khovanskii--Rolle reduction]
\label{prop:rg-relative-KR}
Let $V$ be a smooth $k$-manifold carrying a separating chain
\[
 \Sigma_k\subset\Sigma_{k-1}\subset\cdots\subset\Sigma_0=V
\]
with $\Sigma_r=\Sigma_{r-1}\cap\{\psi_r=0\}$ and
$d\psi_r\ne0$ on $\Sigma_r$.  Let $\rho$ be a covering function on an ambient
manifold whose restriction to every $\Sigma_r$ is covering.  Assume all
intersections have been perturbed to regular levels.  Successively replace the
last remaining separating equation by either its contact equation or, once,
a small covering level.  If $B_0$ denotes the final all-contact count and
$B_j$ the final count when the covering level is first inserted at step $j$,
then
\begin{equation}\label{eq:rg-KR-inequality}
 \#\Sigma_k\le B_0+\frac12\sum_{j=1}^kB_j.
\end{equation}
After a covering level has been inserted, all later reductions have no further
covering term.
\end{proposition}

\begin{proof}
This is the relative version of the one-step argument in Kaloshin's
Khovanskii--Rolle reduction \cite[Theorems~16--18]{Kaloshin2003}.  We recall the point that matters here.  After all
previous rigid/contact levels have been chosen regularly, restrict to a
connected component $C$ of the one-dimensional manifold obtained by imposing
all but the last separating equation.  On $C$ the last separator is a regular
level of a smooth function $f$.  Choose a sufficiently small regular covering
level $\rho=\varepsilon$.  The compact core
$C_\varepsilon=\{\rho\ge\varepsilon\}\cap C$ is a finite union of circles and
compact intervals.

For the circle components, the ordinary Rolle lemma bounds the number of
preimages of the separator level by the number of preimages of a sufficiently
small regular level of the tangential derivative.  On an interval component the
same bound has one additional endpoint term.  Summing over all interval
components gives at most one endpoint contribution per interval, hence one half
of the number of intersections with $\rho=\varepsilon$.  The tangential
derivative vanishes precisely when the current differentials together with the
last Pfaffian form become linearly dependent; relative to any nonvanishing
volume form this is the contact determinant.  Choosing the contact levels
successively from the sufficiently small regular values, exactly as in the
proof of Kaloshin's one-step reduction theorem, therefore gives
\[
 N_{\rm current}\le
 N_{\rm contact}+\frac12N_{\rm cover}.
\]
No nondegeneracy of the zero contact level is assumed.

Iterate this regular-level inequality.  If the covering level is selected at
some step, the new rigid level lies in the compact hypersurface
$\{\rho=\varepsilon\}$, so every subsequent one-dimensional manifold is
compact and the endpoint term disappears.  Thus the induction has one
all-contact branch and exactly one branch for each possible first covering
step, proving \eqref{eq:rg-KR-inequality}.
\end{proof}

For the flight problem take
\begin{equation}\label{eq:rg-S}
 S=[-2M,2M]^k_q\times[0,2]^k_y.
\end{equation}
As before,
\begin{equation}\label{eq:rg-support-base}
 \sum_{j=1}^k\Newt(Q_j)\subseteq S,
 \qquad
 \Newt(P_\varepsilon)\subseteq S.
\end{equation}
In the mixed coframe $d\arg q_j$, $d\log y_j$ on nonzero-slope flights and
$dy_j$ on zero-slope flights, tangential differentiation does not enlarge the
adopted box supersupports.

For $N\ge1$ set
\begin{align}
 C_0(N)&=2^{N(N-1)/2},\label{eq:rg-C0}\\
 C_j(N)&=
 2^{\frac{(j-1)(j-2)}2+\frac{(N-j)(N-j-1)}2}
 (2^{j-1}+1)^{N-j},\qquad 1\le j\le N,\label{eq:rg-Cj}\\
 \Theta_N&=C_0(N)+\frac12\sum_{j=1}^{N}C_j(N).
 \label{eq:rg-Theta}
\end{align}

\begin{theorem}[Unconditional graph-structured fixed-itinerary bound]
\label{thm:rg-fixed-itinerary}
For a prescribed cyclic $k$-flight itinerary of a real scalar periodic
piecewise-linear equation whose slope is constant on every visited region and
whose real trigonometric forcing has Fourier degree at most $M\ge1$, the number of isolated
crossing limit cycles realizing that itinerary satisfies
\begin{equation}\label{eq:rg-final-bound}
 \boxed{
 N_X(k,M)\le
 \Theta_k\,k!(4M)^k\mathsf L_{2k}.}
\end{equation}
Here $\mathsf L_n$ is the Lucas sequence.  The estimate is uniform in the
switching levels and coefficients and includes zero slopes.
\end{theorem}

\begin{proof}
First count nondegenerate transition solutions.  Choose an arbitrary finite
collection $Z$.  By \cref{lem:rg-regularization}, shift the matching level to a
common arbitrarily small regular value $b$ so that every point of $Z$ persists
and the physical partial leaves restrict to a separating chain on
$V_b=Q^{-1}(b)$.  The covering function of \cref{lem:rg-covering} restricts
to a covering function on every member of this chain, because these level
sets are closed in \(\mathcal M_k\).

Apply \cref{prop:rg-relative-KR}.  Fix on $\mathcal M_k$ the nonvanishing
volume form given by the mixed coframe used above.  At every regular rigid
level choose the induced volume form compatibly, so that wedging it with the
differentials of all current rigid equations gives the volume form from the
preceding ambient level.  This is the standard compatible-volume normalization
in the Khovanskii reduction.  With this choice the relative contact function is
represented by the ambient determinant itself, rather than merely by that
determinant times an uncontrolled positive function.  In particular, the first
contact function is the coefficient, relative to the fixed mixed volume form,
of
\[
 dQ_1\wedge\cdots\wedge dQ_k\wedge
 d\phi_1\wedge\cdots\wedge d\phi_k.
\]
Hence it is polynomial--Laurent and its Newton polytope is contained in $S$.
The same compatible normalization applies at every later contact step.  If
$H_r$ is the contact function generated at the $r$th all-contact step, the
support-preserving coframe therefore gives recursively
\begin{equation}\label{eq:rg-contact-support}
 \Newt(H_r)\subseteq2^{r-1}S,
 \qquad r=1,\ldots,k.
\end{equation}
Indeed, the original rigid equations always contribute at most
$\sum_j\Newt(Q_j)\subseteq S$, and the previously generated contact equations
contribute their Minkowski sums.  Hence the all-contact dilation product is
$C_0(k)$.

If the covering equation is first inserted at step $j$, replace the smooth
level $\rho=\varepsilon$ by its cleared Laurent equation
$P_\varepsilon=0$.  The clearing factor is strictly positive on the real
ambient manifold, so this is the same compact regular hypersurface and
$dP_\varepsilon$ is a positive multiple of $d\rho$ along it.  We may therefore
use $P_\varepsilon$ as the rigid defining equation in all subsequent contact
steps.  Its polytope is in $S$, and the compatible-volume normalization then
keeps every later contact polynomial--Laurent.  The later dilation factors obey
the same recurrence $a_r=1+\sum_{\ell<r}a_\ell$ with the covering factor reset
to $1$, giving the product $C_j(k)$ in \eqref{eq:rg-Cj}.  Small regular level
perturbations do not enlarge the boxes because all adopted boxes contain the
origin.

Every final algebraic system contains the $k$ shifted flight equations together
with $k$ contact/covering equations.  Bernstein's theorem and mixed-volume
monotonicity therefore give, on the all-contact branch,
\[
 C_0(k)\operatorname{MV}(L_1,\ldots,L_k,S^k),
\]
and on the branch with first covering insertion at step $j$,
\[
 C_j(k)\operatorname{MV}(L_1,\ldots,L_k,S^k).
\]
The flight--phase incidence graph is the even cycle $C_{2k}$.  The
interaction-graph mixed-volume theorem proved above yields
\begin{equation}\label{eq:rg-cycle-MV}
 \operatorname{MV}(L_1,\ldots,L_k,S^k)
 =k!(4M)^k\mathsf L_{2k}.
\end{equation}
Combining \eqref{eq:rg-KR-inequality} with these BKK bounds gives
\eqref{eq:rg-final-bound} for every finite collection $Z$, and therefore for
all nondegenerate transition solutions.

Finally, Bravo and Trinidad-Forte's additive-bias perturbation shows that an
upper bound for simple crossing cycles bounds all isolated crossing cycles \cite[Proposition~17]{BravoTrinidadForte2026}.
The perturbation changes only the zero Fourier coefficients, preserves
constant slopes, and does not enlarge any Newton box.  Transversality preserves
the prescribed itinerary.  This proves \eqref{eq:rg-final-bound} without a
nondegeneracy assumption.
\end{proof}

\begin{corollary}[Graph-structured constant-slope Hilbert bound]
\label{cor:rg-global}
If $\mathcal H_{\rm cs}^{\times}(n,M)$ denotes the maximal number of crossing
limit cycles in the constant-slope subclass with $n+1$ regions and Fourier
degree at most $M\ge1$, then
\begin{equation}\label{eq:rg-global}
 \mathcal H_{\rm cs}^{\times}(n,M)
 \le2Mn\max_{2\le k\le2Mn}
 \left[\Theta_k k!(4M)^k\mathsf L_{2k}\right].
\end{equation}
Following the global bookkeeping in the proof of the general bound of Bravo and Trinidad-Forte, the corresponding total bound is the right-hand side plus $2$; for continuous piecewise vector fields the same bound also covers tangential cycles.
\end{corollary}

\begin{proof}
Use the crossing bound $k\le2Mn$ and the bound of $2Mn$ crossing sequences of
Bravo and Trinidad-Forte \cite[Propositions~12 and~15]{BravoTrinidadForte2026}, then apply \cref{thm:rg-fixed-itinerary} sequence by
sequence.  Their global bookkeeping supplies the final $+2$ term, and their continuous-case perturbation \cite[Proposition~19]{BravoTrinidadForte2026} removes tangencies while preserving the constant-slope subclass.
\end{proof}

\begin{remark}[Comparison with the direct Pfaffian chain]
For fixed $k$, \eqref{eq:rg-final-bound} has the form $C_kM^k$.  Since
$\mathsf L_{2k}=((3+\sqrt5)/2)^k+((3-\sqrt5)/2)^k$ and
$\Theta_k=2^{k(k-1)/2+O(k)}$, its logarithm has leading quadratic term
$k^2/2$.  At the global crossing scale $k\le2Mn$, this becomes
\[
 \log_2 \mathcal H_{\rm cs}^{\times}(n,M)
 \le 2M^2n^2+O\!\left(Mn\log(MnM)\right).
\]
By contrast, the direct encoding of time-dependent slopes by Bravo and
Trinidad-Forte uses a Pfaffian chain of length $4k\le8Mn$, whose general
Khovanskii factor $2^{r(r-1)/2}$ contributes
\[
 32M^2n^2+O(Mn)
\]
to the quadratic exponent.  Thus, on the constant-slope subclass, the toric
flight encoding reduces the leading quadratic coefficient in this part of the
bound by a factor of $16$, in addition to changing the fixed-itinerary Fourier
degree dependence from the general $M^{4k}$ scale to $M^k$.
\end{remark}
